\documentclass[10pt]{article}

\usepackage[a4paper,margin=1.65cm]{geometry}
\usepackage[T2A]{fontenc}
\usepackage[utf8]{inputenc}
\usepackage{amsmath,amssymb}
\usepackage{enumitem}
\usepackage{hyperref}
\usepackage{titlesec}
\usepackage{xcolor}
\usepackage{microtype}
\usepackage{lmodern}
\usepackage{booktabs}
\usepackage{array}

\definecolor{accent}{HTML}{1F4E79}
\definecolor{textgray}{HTML}{333333}
\definecolor{muted}{HTML}{666666}
\definecolor{linegray}{HTML}{D7DEE8}

\pagestyle{empty}
\hypersetup{colorlinks=true,linkcolor=accent,urlcolor=accent}
\color{textgray}
\setlength{\parindent}{0pt}
\setlength{\parskip}{0.36em}
\setlist[itemize]{leftmargin=1.35em,itemsep=0.08em,topsep=0.12em}
\titleformat{\section}{\large\bfseries\color{accent}}{}{0pt}{}[\vspace{-0.25em}{\color{linegray}\titlerule}]
\titlespacing*{\section}{0pt}{0.85em}{0.28em}
\titleformat{\subsection}{\normalsize\bfseries\color{accent}}{}{0pt}{}
\titlespacing*{\subsection}{0pt}{0.55em}{0.12em}

\newcommand{\norm}[1]{\lVert#1\rVert}
\newcommand{\code}[1]{\texttt{\detokenize{#1}}}

\begin{document}

{\color{accent}\Large\bfseries Sampling plan for Kolmogorov inverse problems}\\
{\color{muted}\small Что уже реализовано, как этим пользоваться, что добавить для conditional sampling}

\section{Что уже сделано}

\subsection{Датасет}

Основной датасет для дальнейших inverse experiments --- velocity snapshots Kolmogorov flow:
\[
    u = (u_x,u_y), \qquad [N,2,64,64].
\]

Данные генерируются в \code{kolmogorov_dataset/}. Симуляция идет в vorticity form на
periodic grid. В текущем варианте поле считается на \(256{\times}256\), затем velocity
усредняется до \(64{\times}64\). В \code{.npz} сохраняются raw values, без нормализации.

Что лежит в одном chunk:
\begin{itemize}
  \item \code{images}: velocity fields \([N,2,H,W]\);
  \item \code{trajectory_id}, \code{snapshot_index}, \code{step}: откуда взят snapshot;
  \item \code{viscosity}, \code{drag}, \code{forcing_amp}: параметры симуляции;
  \item \code{split}: train / val / test.
\end{itemize}

\subsection{Обучение prior}

Основной prior --- VP-SDE score model из \code{score_training/}. Модель предсказывает
noise \(\varepsilon\) из зашумленного поля:
\[
    x_t = \mu_t x_0 + \sigma_t \varepsilon,
    \qquad
    \mu_t=\cos(\omega t),
    \qquad
    \sigma_t=\sqrt{1-\mu_t^2}.
\]

Вход модели:
\[
    [x_t,\sin x,\cos x,\sin y,\cos y].
\]
Координатные Fourier channels добавляются как clean conditioning channels: они не шумятся
и не участвуют в loss. UNet --- \code{diffusers.UNet2DModel} с circular padding.

Нормализация считается только по train split:
\[
    x_{\mathrm{norm}} = \frac{x_{\mathrm{raw}}-\mathrm{mean}}{\mathrm{std}}.
\]
\code{mean} и \code{std} сохраняются в checkpoint. Все samplers должны работать в
normalized space. Метрики и картинки нужно считать в raw space.

\subsection{Unconditional sampling}

В \code{score_training/sde.py} уже есть deterministic VP sampler. На каждом шаге:
\[
    \hat x_0 =
    \frac{x_t-\sigma_t\varepsilon_\theta(x_t,t)}{\mu_t},
    \qquad
    x_{t_{\mathrm{next}}} =
    \mu_{t_{\mathrm{next}}}\hat x_0
    + \sigma_{t_{\mathrm{next}}}\varepsilon_\theta(x_t,t).
\]

Conditional samplers используют тот же reverse process. Сначала модель предсказывает
\(\varepsilon_\theta(x_t,t)\), затем из него считается \(\hat x_0\). После этого sampler
добавляет информацию из observation \(y\): либо меняет \(\hat x_0\), либо делает correction
для \(x_t\), либо склеивает known/unknown regions.

\section{Observation operators}

Для сравнения методов будут использоваться четыре фиксированных оператора наблюдения.
Все они применяются к двум velocity channels одинаково.

\begin{itemize}
  \item \textbf{Regular sparse grid.} Известны значения \(u_x,u_y\) в узлах регулярной
        сетки с шагом 4 по обеим координатам:
        \(M_{ij}=1\) для \(i\bmod 4=0,\; j\bmod 4=0\).
        Observation имеет тот же spatial size \(64{\times}64\), но значения вне mask
        не используются.
  \item \textbf{Central box mask.} Известно полное velocity field внутри центрального окна
        \(32{\times}32\):
        \[
            i,j \in \{16,\ldots,47\}.
        \]
        Вне окна поле неизвестно.
  \item \textbf{Average downsample.} Observation получается average pooling:
        \[
            A: [2,64,64]\rightarrow [2,16,16],
        \]
        kernel \(4{\times}4\), stride 4.
  \item \textbf{Periodic Gaussian blur.} Observation получается circular convolution
        с Gaussian kernel, \(\sigma=2\) grid cells. После blur добавляется Gaussian noise
        с фиксированным \(\sigma_y\).
\end{itemize}

\section{Sampling methods}

\subsection{DDNM}

Используется для линейных операторов, где есть \(A^\dagger\). Самый простой первый
case --- mask/inpainting.

На каждом reverse step:
\[
    \hat x_0 =
    \frac{x_t-\sigma_t\varepsilon_\theta(x_t,t)}{\mu_t}.
\]
Затем заменить observed part:
\[
    \hat x_0^{\mathrm{corr}}
    =
    A^\dagger y + (I-A^\dagger A)\hat x_0.
\]
После этого сделать обычный VP reverse step, но уже из \(\hat x_0^{\mathrm{corr}}\).

Для mask оператор \(A(x)=M\odot x\), а \(A^\dagger(y)=M\odot y\). Поэтому correction
упрощается до замены known pixels:
\[
    \hat x_0^{\mathrm{corr}} = y_{\mathrm{known}} + (1-M)\odot \hat x_0.
\]
Для noisy observations вместо жесткой замены лучше использовать soft update:
\[
    \hat x_0^{\mathrm{corr}}
    =
    \hat x_0 + \lambda_t A^\dagger(y-A\hat x_0).
\]

\subsection{DPS}

DPS работает для любого differentiable \(A\). Он не требует \(A^\dagger\), но требует
gradient через модель.

На каждом шаге \(x_t\) должен быть differentiable. Модель дает \(\hat x_0(x_t,t)\),
после чего считается measurement loss:
\[
    L_y=\norm{A(\hat x_0)-y}_2^2.
\]
Gradient по текущему noisy state:
\[
    g=\nabla_{x_t}L_y
\]
используется как correction:
\[
    x_t^{\mathrm{corr}} = x_t - s_t g.
\]
Дальше reverse step идет из corrected state.

Веса модели не обновляются. Optimizer не нужен. Нужно только autograd по \(x_t\).

Параметры:
\begin{itemize}
  \item \code{guidance_scale}: общий масштаб gradient step;
  \item \code{measurement_sigma}: шум observation;
  \item \code{guidance_start}, \code{guidance_end}: диапазон \(t\), где включается guidance;
  \item gradient clipping или normalization для устойчивости.
\end{itemize}

\subsection{RePaint}

RePaint использовать только для mask. Для blur/downsample он не подходит.

Идея: known region каждый раз заменяется observation, зашумленным до текущего уровня \(t\).
Unknown region генерируется моделью.

Known part строится через forward noising observation:
\[
    x_t^{\mathrm{known}} = \mu_t y + \sigma_t \varepsilon.
\]
Unknown part берется из prior step модели. Затем они склеиваются по mask:
\[
    x_t = M\odot x_t^{\mathrm{known}} + (1-M)\odot x_t^{\mathrm{generated}}.
\]
Jumps добавляются поверх этого процесса: sampler иногда переходит назад к более шумному
состоянию и повторяет несколько reverse steps, чтобы согласовать границу между observed
и generated regions.

\subsection{DDRM}

DDRM нужен для линейных операторов, которые удобно диагонализовать. Для нас это blur и
часть downsample/super-resolution experiments. Делать его после DDNM и DPS.

Идея: работать в basis, где \(A\) раскладывается по singular values:
\[
    A = U\Sigma V^\top.
\]
Для каждого spectral component отдельно смешивается prior prediction и measurement.

Для periodic blur удобно использовать Fourier basis. Blur тогда становится умножением
на transfer function \(H(k)\). Correction считается по каждому \(k\) отдельно, с учетом
\(H(k)\), текущего diffusion noise и measurement noise \(\sigma_y\).

\section{Methods vs observations}

\begin{center}
\begin{tabular}{p{0.18\linewidth}p{0.17\linewidth}p{0.16\linewidth}p{0.17\linewidth}p{0.16\linewidth}}
\toprule
\textbf{Method} & \textbf{Sparse grid} & \textbf{Box mask} &
\textbf{Downsample} & \textbf{Blur} \\
\midrule
DDNM &
yes &
yes &
yes &
yes \\
DPS &
yes &
yes &
yes &
yes \\
RePaint &
yes &
yes &
no &
no \\
DDRM &
not primary &
not primary &
yes &
yes \\
\bottomrule
\end{tabular}
\end{center}

DDNM and RePaint are most natural for masks. DPS is the general differentiable
baseline and can be run for all four operators. DDRM is reserved for operators with
a convenient spectral representation: downsample and blur.

\section{Physics-informed sampling}

Для Kolmogorov velocity fields основное физическое ограничение:
\[
    \nabla\cdot u = 0.
\]
Его можно добавлять к sampling без переобучения модели. Базовый вариант --- Helmholtz
projection в Fourier space:
\[
    \hat u_{\mathrm{proj}}(k)
    =
    \hat u(k)
    -
    k\,\frac{k\cdot \hat u(k)}{\norm{k}^2},
    \qquad k\ne0.
\]
Она удаляет compressible component и оставляет divergence-free part поля. Projection
нужно применять в normalized model space или в raw space согласованно с тем, где делается
correction. Для финальных метрик поле все равно переводится в raw space.

\subsection{DDNM + physics}

В DDNM physics удобно применять к \(\hat x_0^{\mathrm{corr}}\). Сначала делается
data-consistency correction:
\[
    \hat x_0^{\mathrm{corr}}
    =
    A^\dagger y + (I-A^\dagger A)\hat x_0,
\]
после этого к corrected estimate применяется Helmholtz projection:
\[
    \hat x_0^{\mathrm{phys}} = P_{\mathrm{divfree}}(\hat x_0^{\mathrm{corr}}).
\]
Дальше VP reverse step использует \(\hat x_0^{\mathrm{phys}}\).

Для mask observations есть важный trade-off: hard projection может немного изменить known
pixels. Если нужно строго сохранить observation, можно делать projection только на unknown
part или после projection повторно вставлять known pixels.

\subsection{DPS + physics}

В DPS physics можно добавить как дополнительный guidance term:
\[
    L = L_y + \lambda_{\mathrm{div}} L_{\mathrm{div}},
    \qquad
    L_{\mathrm{div}} = \norm{\nabla\cdot \hat x_0}_2^2.
\]
Тогда gradient correction считается по общей функции \(L\):
\[
    x_t^{\mathrm{corr}}
    =
    x_t - s_t \nabla_{x_t}L.
\]
Это soft physics guidance. Его удобно использовать вместе с differentiable operators,
потому что DPS уже требует autograd по \(x_t\).

Альтернативный вариант --- не добавлять \(L_{\mathrm{div}}\), а после DPS correction
проецировать \(\hat x_0\) через Helmholtz projection. Это проще и стабильнее, но может
конфликтовать с observation consistency для mask.

\subsection{RePaint + physics}

В RePaint known region задается observation, unknown region генерируется prior-ом. Поэтому
physics projection лучше применять к generated/unknown part, а known part затем снова
вставлять по mask:
\[
    x = M\odot x_{\mathrm{known}} + (1-M)\odot P_{\mathrm{divfree}}(x_{\mathrm{generated}}).
\]
Если применить projection ко всему полю сразу, known pixels могут измениться. Это плохо для
inpainting setting, где observed region должен оставаться фиксированным.

\subsection{DDRM + physics}

В DDRM data update происходит в spectral/SVD basis. Physics projection тоже удобно делать
в Fourier space, поэтому для blur этот вариант естественный: после spectral DDRM update
можно применить Helmholtz projection к velocity field.

Для downsample projection тоже возможна, но порядок важен. Практичный вариант:
сначала выполнить DDRM measurement update, затем Helmholtz projection, затем следующий
reverse step. Если projection заметно ухудшает consistency с low-resolution observation,
можно добавить повторную мягкую correction по \(A(\hat x_0)-y\).

\section{Evaluation pipeline}

Evaluation должен быть одинаковым для всех методов. Для held-out field
\(u_{\mathrm{true}}\) строится observation
\[
    y = A(u_{\mathrm{true}}) + \eta.
\]
Sampler получает один и тот же checkpoint, один и тот же operator \(A\), один и тот же
noise level и фиксированный seed. Внутри sampler все вычисления идут в normalized space.
После sampling reconstruction переводится обратно в raw velocity space, и уже там
считаются метрики и строятся plots.

Для каждого запуска сохраняется одна строка в CSV. Для небольшой подвыборки сохраняются
PNG visualizations: ground truth, observation, reconstruction, error map и vorticity.

Минимальный набор experiments:
\begin{itemize}
  \item regular sparse grid, stride 4;
  \item central box mask, \(32{\times}32\);
  \item downsample: \(64{\times}64 \rightarrow 16{\times}16\);
  \item Gaussian blur, \(\sigma=2\), with Gaussian observation noise.
\end{itemize}

Для debug run достаточно \(N=16\), one seed. Для финальной таблицы: \(N=256\), seeds
\(\{0,1,2\}\).

\section{Metrics}

Все метрики считаются в raw velocity space.

\begin{itemize}
  \item Relative L2:
        \[
            \frac{\norm{\hat u-u_{\mathrm{true}}}_2}
                 {\norm{u_{\mathrm{true}}}_2}.
        \]
  \item RMSE:
        \[
            \sqrt{\frac{1}{2HW}\norm{\hat u-u_{\mathrm{true}}}_2^2}.
        \]
  \item Measurement consistency:
        \[
            \frac{\norm{A(\hat u)-y}_2}{\norm{y}_2}.
        \]
  \item Divergence:
        \[
            \frac{\norm{\nabla\cdot \hat u}_2}{\norm{\hat u}_2}.
        \]
  \item Vorticity RMSE. Сначала считать
        \[
            \omega = \partial_x u_y - \partial_y u_x,
        \]
        затем RMSE между \(\omega(\hat u)\) и \(\omega(u_{\mathrm{true}})\).
\end{itemize}

CSV columns:
\begin{verbatim}
sample_id, seed, method, operator, noise_sigma,
rel_l2, rmse, measurement_error, divergence, vorticity_rmse
\end{verbatim}

\end{document}
